% Exercise 3.14

\begin{itemize}
    \item[(a)] The form of the linear model is given by
        \begin{equation*}
            \begin{split}
                Y &= 2 + 2X_1 + 0.3X_2 + \varepsilon_1 \\
                  &= 2 + 2X_1 + 0.3\biggl(0.5X_1 + \frac{\varepsilon_2}{10}\biggr)
                      + \varepsilon_1 \\
                  &= 2 + 2X_1 + 0.15X_1 + 0.03\varepsilon_2 + \varepsilon_1 \\
                  &= 2 + 2.15X_1 + \varepsilon_3,
            \end{split}
        \end{equation*}
        for some $\varepsilon_3 = \varepsilon_1 + 0.03\varepsilon_2$, which is simply
        an error term. The regression coefficients are $\beta_0 = 2$, $\beta_1 = 2$,
        and $\beta_2 = 0.03$.
    \item[(b)] From (a), we see that $X_2 = 0.5X_1 + \varepsilon_2/10$, so that the
        variables show some correlation. In Figure~\ref{fig3_14scat}, we see a 
        scatterplot displaying the relationship between the variables.
        \begin{figure}[!ht]
            \includegraphics[width=\textwidth, center]{../plots/ex3_14_b.pdf}
            \caption{Scatterplot of $X_1$ and $X_2$ \label{fig3_14scat}}
        \end{figure}
    \item[(c)] Using the given model, we have the following results for the least
        squares regression prediction of \verb|y| using \verb|x1| and \verb|x2|:
        \scriptsize\begin{verbatim}
> lm.fit <- lm(y ~ x1 + x2)
> summary(lm.fit)

Call:
lm(formula = y ~ x1 + x2)

Residuals:
    Min      1Q  Median      3Q     Max 
-2.8311 -0.7273 -0.0537  0.6338  2.3359 

Coefficients:
            Estimate Std. Error t value Pr(>|t|)    
(Intercept)   2.1305     0.2319   9.188 7.61e-15 ***
x1            1.4396     0.7212   1.996   0.0487 *  
x2            1.0097     1.1337   0.891   0.3754    
---
Signif. codes:  0 ‘***’ 0.001 ‘**’ 0.01 ‘*’ 0.05 ‘.’ 0.1 ‘ ’ 1

Residual standard error: 1.056 on 97 degrees of freedom
Multiple R-squared:  0.2088,    Adjusted R-squared:  0.1925 
F-statistic:  12.8 on 2 and 97 DF,  p-value: 1.164e-05
        \end{verbatim}\normalsize
        The regression predicts values of $\hat{\beta}_0 = 2.1305$, 
        $\hat{\beta}_1 = 1.4396$, and $\hat{\beta}_2 = 1.0097$, but the $p$-values for
        each variable are given by $7.61\times {10}^{-15}$, 0.0487, and 0.3754,
        respectively. Thus it would not be fair to reject the null hypothesis 
        $H_0 : \beta_1 = 0$ or $H_0 : \beta_2 = 0$, since the least squares regression
        does not give us enough certainty about the estimates to conclude anything for
        these two coefficients.
    \item[(d)] We have the following least squares regression for predicting \verb|y| 
        from \verb|x1| alone:
        \scriptsize\begin{verbatim}
> lm.fit <- lm(y ~ x1)
> summary(lm.fit)

Call:
lm(formula = y ~ x1)

Residuals:
     Min       1Q   Median       3Q      Max 
-2.89495 -0.66874 -0.07785  0.59221  2.45560 

Coefficients:
            Estimate Std. Error t value Pr(>|t|)    
(Intercept)   2.1124     0.2307   9.155 8.27e-15 ***
x1            1.9759     0.3963   4.986 2.66e-06 ***
---
Signif. codes:  0 ‘***’ 0.001 ‘**’ 0.01 ‘*’ 0.05 ‘.’ 0.1 ‘ ’ 1

Residual standard error: 1.055 on 98 degrees of freedom
Multiple R-squared:  0.2024,    Adjusted R-squared:  0.1942 
F-statistic: 24.86 on 1 and 98 DF,  p-value: 2.661e-06
        \end{verbatim}\normalsize
        In this case, we have much a smaller $p$-value of $2.66 \times {10}^{-6}$ 
        for the coefficient estimate $\hat{\beta}_1 = 1.9759$, so we may reject the 
        null hypothesis $H_0 : \beta_1 = 0$.
    \item[(e)]We have the following least squares regression for predicting \verb|y| 
        from \verb|x2| alone:
        \scriptsize\begin{verbatim}
> lm.fit <- lm(y ~ x2)
> summary(lm.fit)

Call:
lm(formula = y ~ x2)

Residuals:
     Min       1Q   Median       3Q      Max 
-2.62687 -0.75156 -0.03598  0.72383  2.44890 

Coefficients:
            Estimate Std. Error t value Pr(>|t|)    
(Intercept)   2.3899     0.1949   12.26  < 2e-16 ***
x2            2.8996     0.6330    4.58 1.37e-05 ***
---
Signif. codes:  0 ‘***’ 0.001 ‘**’ 0.01 ‘*’ 0.05 ‘.’ 0.1 ‘ ’ 1

Residual standard error: 1.072 on 98 degrees of freedom
Multiple R-squared:  0.1763,    Adjusted R-squared:  0.1679 
F-statistic: 20.98 on 1 and 98 DF,  p-value: 1.366e-05
        \end{verbatim}\normalsize
        In this case, we also have much a smaller $p$-value of $1.37 \times {10}^{-5}$ 
        for the coefficient estimate $\hat{\beta}_2 = 2.8996$, so we may reject the null 
        hypothesis $H_0 : \beta_2 = 0$.
    \item[(f)] No, these results do not contradict eachother, since we have two collinear
        variables, so measuring the individual effect of either \verb|x1| or \verb|x2| on
        \verb|y| alone will not yield conclusive results.
    \item[(g)] In Figures~\ref{fig3_14resi} to~\ref{fig3_14resiii}, we have diagnostic 
        plots for each of the three models in (c)-(e). We can see clearly in 
        Figure~\ref{fig3_14resi} that the new observation is not an outlier, but is
        clearly a high-leverage point.

        In Figure~~\ref{fig3_14resii}, it is an outlier but not a high-leverage point, 
        and in Figure~\ref{fig3_14resiii} it is a high-leverage point but not an outlier.
        
        These judgements are a reflection of the location of the 101st observation in 
        each of the diagnostic plots.

        \begin{figure}[!ht]
            \includegraphics[width=\textwidth, center]{../plots/ex3_14_g_i.pdf}
            \caption{Diagnostic plots for model from  3.14.(c) \label{fig3_14resi}}
        \end{figure}
        \begin{figure}[!ht]
            \includegraphics[width=\textwidth, center]{../plots/ex3_14_g_ii.pdf}
            \caption{Diagnostic plots for model from  3.14.(d) \label{fig3_14resii}}
        \end{figure}
        \begin{figure}[!ht]
            \includegraphics[width=\textwidth, center]{../plots/ex3_14_g_iii.pdf}
            \caption{Diagnostic plots for model from  3.14.(e) \label{fig3_14resiii}}
        \end{figure}
\end{itemize}
